% !TEX root = ../Main.tex
\chapter{标量和矢量}

\section{为什么要有方向}

在 $xOy$ 平面里描述一个点的位置，先看一个对比。

只告诉你"该点到原点 $O$ 的距离为 $L$"，但不告诉方向——满足条件的点有\textbf{无数个}，它们排成一个半径为 $L$ 的圆。可一旦补上方向，比如"在 $x$ 轴正半轴上、到 $O$ 的距离为 $L$"，这个点的位置就\textbf{唯一}确定了。

\begin{center}
\begin{tikzpicture}[>=Stealth, scale=1]
    \draw[->] (-1.6,0) -- (1.9,0) node[right] {$x$};
    \draw[->] (0,-1.6) -- (0,1.6) node[above] {$y$};
    \draw[blue] (0,0) circle (1.2);
    \filldraw[red] (1.2,0) circle (1.5pt) node[above right] {\footnotesize $(L,0)$};
\end{tikzpicture}
\end{center}

可见，光有"大小"$L$ 不够，还得有"方向"才能定下位置——物理量正是因此分成两类。

\section{标量与矢量}

\state{标量与矢量}{
    \textbf{标量}：只有大小。\quad 例如路程 $L$。\\[2pt]
    \textbf{矢量}：既有大小，也有方向。\quad 例如位移 $S$。
}

路程 $L$ 与位移 $S$ 正好是这对概念的代表：

\[
\text{路程 }L\ \xrightarrow{\quad\text{加上方向}\quad}\ \text{位移 }S.
\]

\textbf{位移}是由\textbf{初位置}指向\textbf{末位置}的有向线段——它只看起点和终点，不管中途绕了多远。

\begin{center}
\begin{tikzpicture}[>=Stealth, scale=1]
    \draw[blue, thick] (0,0) .. controls (0.6,1.0) and (1.8,-0.6) .. (2.6,0.4);
    \filldraw (0,0) circle (1.3pt) node[below left] {$A$};
    \filldraw (2.6,0.4) circle (1.3pt) node[above right] {$B$};
    \draw[->, very thick, red] (0,0) -- (2.6,0.4) node[midway, below right] {$S$};
\end{tikzpicture}
\end{center}

\section{矢量的变化与比较}

\textbf{矢量的变化}：大小和方向，只要其中\textbf{有一个}发生变化，就认为这个矢量变了。

\textbf{矢量的比较}：两个矢量\textbf{相等（相同）}，要求大小和方向\textbf{都一致}。至于"哪个大哪个小"，比较的只是它们的大小。

\rmk{这条在后面非常有用：匀速圆周运动里速度大小不变、方向一直在变，所以速度这个矢量一直在变。}
